Read by its binary digits rather than by its magnitude,
a non-negative integer $b$ is a finite set of positions:
bit $k$ of $b$ is set exactly when $k$ belongs to the set.
In the book of Section \ref{sec.fasterBook}
the position of a price is its distance from an edge of the band,
and one integer is the set of prices occupied on one side.
Table \ref{tab.bitsAsSets} lists the operations of Python that act on the integer as on the set.

\begin{table}[h!]
 \centering
 \begin{tabular}{@{}ll@{}}
  \toprule
  Python & on the set \\
  \midrule
  \codehl{1 << k} & the singleton $\lbrace k \rbrace$ \\
  \codehl{b | (1 << k)} & $k$ added \\
  \codehl{b \& \textasciitilde(1 << k)} & $k$ removed \\
  \codehl{b >> k \& 1} & $1$ if $k$ belongs to it, $0$ otherwise \\
  \codehl{b.bit\_count()} & the number of its elements \\
  \bottomrule
 \end{tabular}
 \caption{An integer as a set of positions.}
 \label{tab.bitsAsSets}
\end{table}

The integer $360 = 2^{8} + 2^{6} + 2^{5} + 2^{3}$ is the set $\lbrace 3, 5, 6, 8 \rbrace$,
and Listing \ref{listing.bitsAsSets} builds it and edits it.
Python's integers are unbounded, so the set is as large as the band;
where an integer is a machine word, a wider set is an array of words,
with the arithmetic of the indices written out.

\begin{lstlisting}[caption={The integer $360$ as the set $\lbrace 3, 5, 6, 8 \rbrace$.}, label=listing.bitsAsSets]
b = 0
for k in (3, 5, 6, 8):
    b |= 1 << k
assert b == 360
assert [k for k in range(10) if b >> k & 1] == [3, 5, 6, 8]
assert b.bit_count() == 4
b &= ~(1 << 5)
assert [k for k in range(10) if b >> k & 1] == [3, 6, 8]
\end{lstlisting}
